\chapter{Continuity and the Intermediate Value Theorem}

There is a word we use loosely in everyday life and precisely in mathematics:
\emph{continuous}. We say a road runs continuously from one town to the next, meaning
you can drive the whole way without ever being teleported across a gap. We say someone
worked continuously for three hours, meaning no interruptions. Underneath both uses is
the same picture: \emph{no sudden jumps}. You can get from here to there by moving a
little at a time.

Calculus is built on this picture. Almost every theorem you will meet later --- that a
function must hit every value between two of its outputs, that it must attain a highest
and a lowest point on a closed interval, that the area under a curve makes sense, that a
differentiable function is well behaved --- rests on continuity. So before we can
differentiate or integrate anything with confidence, we need to say, with total
precision, what it means for a function to have ``no jumps.''

The remarkable thing is that the right definition does not mention jumps at all. It is
phrased entirely in the language of limits you built in the previous chapters. By the end
of this chapter you will see \emph{why} the limit definition captures exactly the
no-jump idea, you will be able to classify every way a function can fail to be
continuous, and you will prove and use two of the great existence theorems of analysis:
the Intermediate Value Theorem and the Extreme Value Theorem.

\section{From ``no jumps'' to a definition}
\label{ch08:no-jumps}

Let us start with the picture and refuse to leave it until the definition feels forced
upon us.

Imagine dragging a pencil along the graph of a function $f$. Continuity at a point $a$
should mean: as your pencil approaches the vertical line $x=a$ from either side, the
height of the curve approaches the height $f(a)$ of the point actually sitting above $a$.
Nothing jumps; nothing is missing; the point you would \emph{predict} by sliding in is
exactly the point that is \emph{there}.

That last sentence is the whole idea, and it splits into three separate promises, each of
which can fail independently.

\begin{intuitionbox}
For $f$ to be continuous at $a$, three things must all hold:
\begin{enumerate}
  \item $f(a)$ must \emph{exist} --- there has to be a point on the graph above $a$.
  \item $\displaystyle\lim_{x\to a} f(x)$ must \emph{exist} --- the curve must approach a
        single definite height as you slide in from both sides.
  \item These two numbers must \emph{agree} --- the height you slide toward equals the
        height that is actually there.
\end{enumerate}
The first is about the point. The second is about the approach. The third demands that
the approach and the point match. Continuity is precisely the demand that prediction and
reality coincide.
\end{intuitionbox}

Compressing all three promises into one line gives the definition.

\begin{definition}[Continuity at a point]
\label{ch08:def-continuity}
A function $f$ is \textbf{continuous at} $a$ if
\[
  \lim_{x\to a} f(x) = f(a).
\]
We mean this in the strongest sense: $a$ is in the domain of $f$ (so $f(a)$ is a real
number), the limit on the left exists, and the two are equal. If $f$ is continuous at
every point of a set $S$, we say $f$ is \textbf{continuous on} $S$.
\end{definition}

This single equation quietly bundles all three promises. If $f(a)$ did not exist, the
right side would be meaningless. If the limit did not exist, the left side would be
meaningless. And if both existed but disagreed, the equation would simply be false. The
equation $\lim_{x\to a} f(x) = f(a)$ asserts all three at once. That economy is exactly
why mathematicians chose it.

\begin{happeningbox}
Why does an equation about limits capture a \emph{geometric} idea about jumps? Because a
limit \emph{is} the machinery of ``sliding in.'' Recall what $\lim_{x\to a} f(x)=L$ really
asserts: for every tolerance $\eps>0$, no matter how small, there is a window of inputs
$0<\abs{x-a}<\delta$ on which $f$ stays within $\eps$ of $L$. Continuity says we may take
$L=f(a)$ \emph{and drop the restriction} $x\ne a$. In one breath:
\[
  \text{for every } \eps>0 \text{ there is } \delta>0 \text{ such that }
  \abs{x-a}<\delta \implies \abs{f(x)-f(a)}<\eps.
\]
Read that aloud as a promise about jumps: ``if you want my output to land within $\eps$
of $f(a)$, just keep your input within $\delta$ of $a$.'' A jump is exactly a failure of
this promise --- a height the output refuses to approach no matter how close the input
gets. Continuity is the guarantee that small input changes force small output changes.
There is no gap to teleport across.
\end{happeningbox}

\begin{focusbox}
Notice the subtle but crucial difference from the bare limit definition of the previous
chapter. There we wrote $0<\abs{x-a}<\delta$, deliberately \emph{excluding} the point
$x=a$, because a limit is about the neighborhood, not the point itself. Continuity is the
stronger statement that we may \emph{include} $x=a$: the inequality $\abs{x-a}<\delta$ now
allows $x=a$, where it reads $\abs{f(a)-f(a)}=0<\eps$ for free. The limit ignores the
point; continuity insists the point cooperate with its neighbors.
\end{focusbox}

\subsection{One-sided continuity and continuity on an interval}

Sometimes a function only ``slides in'' nicely from one side --- think of $\sqrt{x}$ at
$x=0$, where there are no inputs to the left at all. We handle this with one-sided
limits.

\begin{definition}[One-sided continuity]
\label{ch08:def-onesided}
$f$ is \textbf{continuous from the right} at $a$ if $\lim_{x\to a^+} f(x)=f(a)$, and
\textbf{continuous from the left} at $a$ if $\lim_{x\to a^-} f(x)=f(a)$. A function is
continuous at an interior point exactly when it is continuous from both sides there.
\end{definition}

This lets us say cleanly what it means for $f$ to be continuous on a \emph{closed}
interval $[a,b]$: continuous at every interior point, continuous from the right at the
left endpoint $a$, and continuous from the left at the right endpoint $b$. We only ever
ask a function to behave with respect to inputs that actually exist. This convention will
matter enormously when we state the Extreme Value Theorem, which needs a closed interval.

\section{A gallery of discontinuities}
\label{ch08:gallery}

The fastest way to understand continuity is to study how it breaks. Each of the three
promises can fail, and the failures look different on a graph. Let us walk through the
zoo.

\subsection{Removable discontinuity: the patchable hole}

Here the limit exists --- the curve clearly aims at a single height --- but either $f(a)$
is undefined or it is defined ``in the wrong place.'' The graph has a single missing or
misplaced point. The defect is \emph{removable}: redefining $f$ at the one point $a$ to
equal the limit repairs continuity completely.

\begin{figure}[H]
\centering
\begin{tikzpicture}
\begin{axis}[
  width=9.5cm, height=6cm,
  axis lines=middle,
  xlabel={$x$}, ylabel={$y$},
  xmin=-0.5, xmax=4.5, ymin=-0.5, ymax=4.5,
  xtick={2}, xticklabels={$a$}, ytick={\empty},
  samples=100, domain=-0.5:4.5,
  clip=false,
]
\addplot[curveblue, very thick, domain=-0.3:1.97]{0.6*x+1};
\addplot[curveblue, very thick, domain=2.03:4.3]{0.6*x+1};
\addplot[only marks, mark=o, mark size=3pt, curveblue, thick] coordinates {(2,2.2)};
\addplot[only marks, mark=*, mark size=2.5pt, trapred] coordinates {(2,3.6)};
\node[trapred, anchor=west] at (axis cs:2.15,3.6) {\small $f(a)$};
\node[curveblue, anchor=east] at (axis cs:1.85,2.05) {\small $\lim_{x\to a}f$};
\end{axis}
\end{tikzpicture}
\caption{A removable discontinuity. The curve aims at the open circle, but the actual
value $f(a)$ (filled dot) sits elsewhere. The limit exists; it just disagrees with
$f(a)$. Move the dot down into the hole and continuity is restored.}
\label{ch08:fig-removable}
\end{figure}

\begin{example}[The classic hole]
\label{ch08:ex-hole}
Consider
\[
  g(x)=\frac{x^2-1}{x-1}.
\]
At $x=1$ the formula reads $0/0$, undefined, so promise (1) fails: $g(1)$ does not exist.
But for every $x\ne 1$ we may cancel:
\[
  g(x)=\frac{(x-1)(x+1)}{x-1}=x+1,
\]
so $\lim_{x\to 1} g(x)=\lim_{x\to 1}(x+1)=2$. The limit exists and equals $2$; only the
point is missing. Defining $g(1)=2$ plugs the hole and makes $g$ continuous everywhere.
\end{example}

\begin{trapbox}
A removable discontinuity is still a genuine discontinuity. The function $g$ above is
\emph{not} continuous at $1$ as originally written --- it has no value there, so the
equation $\lim_{x\to 1}g(x)=g(1)$ cannot hold. People sometimes say ``$g$ is basically
$x+1$.'' It \emph{becomes} $x+1$ only after you supply the missing value. Until you do,
there is a hole, and a hole is a discontinuity. Cancellation in a formula does not by
itself mend the function; it tells you which value would mend it.
\end{trapbox}

\subsection{Jump discontinuity: the cliff}

Here both one-sided limits exist but \emph{disagree}. The curve approaches one height from
the left and a different height from the right. There is a genuine cliff. No single
redefinition of $f(a)$ can fix this, because there are two different heights competing and
$f(a)$ can only be one number.

\begin{figure}[H]
\centering
\begin{tikzpicture}
\begin{axis}[
  width=9.5cm, height=6cm,
  axis lines=middle,
  xlabel={$x$}, ylabel={$y$},
  xmin=-0.5, xmax=4.5, ymin=-0.5, ymax=4.5,
  xtick={2}, xticklabels={$a$}, ytick={\empty},
  clip=false,
]
\addplot[curveblue, very thick, domain=-0.3:2]{0.4*x+0.6};
\addplot[curveblue, very thick, domain=2:4.3]{0.4*x+1.7};
\addplot[only marks, mark=*, mark size=2.5pt, curveblue] coordinates {(2,1.4)};
\addplot[only marks, mark=o, mark size=3pt, curveblue, thick] coordinates {(2,2.5)};
\node[curveblue, anchor=east] at (axis cs:1.9,1.0) {\small left limit};
\node[curveblue, anchor=west] at (axis cs:2.1,2.9) {\small right limit};
\end{axis}
\end{tikzpicture}
\caption{A jump discontinuity. The left-hand limit and right-hand limit both exist but
are unequal, so the two-sided limit does not exist. The gap between the dots is the size
of the jump.}
\label{ch08:fig-jump}
\end{figure}

\begin{example}[A step function]
\label{ch08:ex-step}
Let $H(x)=0$ for $x<0$ and $H(x)=1$ for $x\ge 0$ (the Heaviside step, the idealized
``switch turns on'' function). Then $\lim_{x\to 0^-}H(x)=0$ while $\lim_{x\to 0^+}H(x)=1$.
The one-sided limits disagree, so $\lim_{x\to 0}H(x)$ does not exist, and $H$ is
discontinuous at $0$. It \emph{is} continuous from the right there, since
$\lim_{x\to 0^+}H(x)=1=H(0)$.
\end{example}

\begin{applicationbox}
Jump discontinuities are not pathological --- they are everywhere in modeling. A light
switch, a digital signal flipping from $0$ to $1$, the price of a stock the instant news
breaks, the charge on a capacitor the moment a switch closes: all are modeled by jumps.
This is exactly why engineers love the Heaviside function and its cousin, the Dirac delta.
Calculus copes with such functions by working on the pieces \emph{between} jumps, where
everything is continuous, and treating the jumps as boundaries. ``Piecewise continuous''
is the honest description of most real-world signals.
\end{applicationbox}

\subsection{Infinite discontinuity: the pole}

Here a one-sided limit (or both) is infinite. The curve shoots off to $\pm\infty$ near
$a$. There is a vertical asymptote. No finite value of $f(a)$ could possibly match an
infinite approach.

\begin{figure}[H]
\centering
\begin{tikzpicture}
\begin{axis}[
  width=9.5cm, height=6cm,
  axis lines=middle,
  xlabel={$x$}, ylabel={$y$},
  xmin=-3.2, xmax=3.2, ymin=-4.2, ymax=4.2,
  xtick={0}, ytick={\empty},
  samples=200,
  restrict y to domain=-12:12,
]
\addplot[curveblue, very thick, domain=-3.1:-0.27]{1/x};
\addplot[curveblue, very thick, domain=0.27:3.1]{1/x};
\draw[trapred, dashed, thick] (axis cs:0,-4.2)--(axis cs:0,4.2);
\end{axis}
\end{tikzpicture}
\caption{An infinite discontinuity: $f(x)=1/x$ near $x=0$. The left limit is $-\infty$,
the right limit is $+\infty$. The dashed line is the vertical asymptote.}
\label{ch08:fig-infinite}
\end{figure}

\subsection{Oscillatory discontinuity: the blur}

The subtlest failure: the function neither settles down nor blows up; it oscillates
forever, infinitely fast, as $x\to a$. The limit fails not because of a jump or a pole but
because there is no single height the values approach at all.

\begin{example}[Endless wiggling]
\label{ch08:ex-osc}
Let $s(x)=\sin(1/x)$ for $x\ne 0$. As $x\to 0$, the argument $1/x$ races off to infinity,
so $\sin(1/x)$ sweeps through every value in $[-1,1]$ infinitely many times in any
neighborhood of $0$. There is no limiting height. No definition of $s(0)$ can rescue
continuity, because the values do not approach \emph{anything}.

Contrast its tamer relative $x\sin(1/x)$: the wild oscillation is throttled by the factor
$x\to 0$, the wiggles get squeezed to zero amplitude, and $\lim_{x\to 0}x\sin(1/x)=0$ by
the Squeeze Theorem. Defining the value at $0$ to be $0$ makes \emph{that} one continuous.
The difference between a curable and an incurable wiggle is whether the amplitude dies
out.
\end{example}

\begin{figure}[H]
\centering
\begin{tikzpicture}
\begin{axis}[
  width=10.5cm, height=5.4cm,
  axis lines=middle,
  xlabel={$x$}, ylabel={$y$},
  xmin=-0.35, xmax=0.35, ymin=-1.25, ymax=1.25,
  xtick={0}, ytick={-1,1},
  samples=200, domain=-0.33:0.33,
  restrict y to domain=-1.2:1.2,
]
\addplot[curveblue, thick]{sin(deg(1/x))};
\end{axis}
\end{tikzpicture}
\caption{$\sin(1/x)$ near the origin. The oscillation never slows; the values fill the
band $[-1,1]$ densely as $x\to 0$, so no limit exists.}
\label{ch08:fig-osc}
\end{figure}

\begin{walkbox}
Take a walk and sort the discontinuities by ``how badly broken'' they are. The removable
hole is barely broken --- one point out of place. The jump is broken but tame --- two
honest one-sided limits. The pole is broken and unbounded. The oscillation is broken and
not even \emph{trying} to have a limit. Now ask yourself: which of these can a computer
plotting $200$ points possibly miss? (Hint: think about whether your sample points might
land on the ``wrong'' side of the trouble, or skip over a thin spike. This is why we prove
things instead of trusting pictures.)
\end{walkbox}

\section{Building continuous functions}
\label{ch08:building}

If we had to check the limit definition by hand for every function, calculus would be
unbearable. Fortunately continuity is \emph{hereditary}: continuous functions combine to
make more continuous functions. We get to know a few atoms are continuous, then assemble
everything else for free.

\subsection{The atoms}

The following are continuous everywhere on their natural domains; each follows directly
from the limit laws of the previous chapters.

\begin{itemize}
  \item Constant functions $f(x)=c$ and the identity $f(x)=x$. (For the identity,
        $\lim_{x\to a}x=a=f(a)$ --- this is almost the definition of a limit.)
  \item Polynomials, everywhere on $\R$.
  \item Rational functions $p(x)/q(x)$, everywhere $q(x)\ne 0$.
  \item $\sin x$, $\cos x$, everywhere; $\exp x$ and, on $(0,\infty)$, $\ln x$.
  \item $\sqrt{x}$ on $[0,\infty)$ (continuous from the right at $0$), and more
        generally $x^r$ on its domain.
\end{itemize}

\subsection{The combination laws}

\begin{theorem}[Continuity of combinations]
\label{ch08:thm-combinations}
Suppose $f$ and $g$ are continuous at $a$. Then so are
\[
  f+g,\qquad f-g,\qquad fg,\qquad cf \ (c\in\R),
\]
and $f/g$ \emph{provided} $g(a)\ne 0$.
\end{theorem}

\begin{whybox}
Why is this true with no new work at all? Because each statement is just a limit law
rephrased. Continuity of $f$ at $a$ means $\lim_{x\to a}f(x)=f(a)$, and likewise for $g$.
The sum law for limits says
\[
  \lim_{x\to a}\bigl(f(x)+g(x)\bigr)=\lim_{x\to a}f(x)+\lim_{x\to a}g(x)=f(a)+g(a)=(f+g)(a),
\]
which is exactly the statement that $f+g$ is continuous at $a$. The product, difference,
and quotient cases are the corresponding limit laws, read the same way. The proviso
$g(a)\ne 0$ in the quotient is precisely the proviso in the limit quotient law: you may
not divide by a quantity heading to zero. \emph{Every} continuity combination law is a
limit law wearing a different hat.
\end{whybox}

This single observation explains the whole list of atoms. A polynomial is built from the
identity and constants by multiplication and addition; apply Theorem~\ref{ch08:thm-combinations}
finitely many times and continuity is guaranteed at every point. A rational function is a
quotient of polynomials, continuous wherever the denominator is nonzero. You never have to
return to $\eps$ and $\delta$ for these.

\subsection{Composition: the most important law}

The deepest of the combination laws concerns \emph{chaining} functions, $g$ then $f$. It
will reappear as the engine of the chain rule in differentiation, so it is worth seeing
clearly now.

\begin{theorem}[Continuity of compositions]
\label{ch08:thm-composition}
If $g$ is continuous at $a$ and $f$ is continuous at $g(a)$, then the composition
$f\circ g$, defined by $(f\circ g)(x)=f(g(x))$, is continuous at $a$.
\end{theorem}

\begin{intuitionbox}
Picture two machines in a row. The first, $g$, takes $x$ near $a$ and (being continuous)
produces output near $g(a)$. The second, $f$, takes inputs near $g(a)$ and (being
continuous there) produces outputs near $f(g(a))$. Feed a small wiggle in $x$ through the
pair: $g$ converts it to a small wiggle near $g(a)$, and $f$ converts that to a small
wiggle near $f(g(a))$. Small in, small out, all the way through. The chain of two no-jump
machines is itself a no-jump machine.
\end{intuitionbox}

\begin{proof}
We use the $\eps$--$\delta$ form of continuity. Let $\eps>0$ be given; we must find
$\delta>0$ so that $\abs{x-a}<\delta$ forces $\abs{f(g(x))-f(g(a))}<\eps$.

Work from the outside in. Because $f$ is continuous at the point $b:=g(a)$, there is an
$\eta>0$ (the tolerance $f$ needs on \emph{its} input) such that
\[
  \abs{y-b}<\eta \implies \abs{f(y)-f(b)}<\eps.
\]
Now hand that $\eta$ to $g$. Because $g$ is continuous at $a$, there is a $\delta>0$ such
that
\[
  \abs{x-a}<\delta \implies \abs{g(x)-g(a)}<\eta.
\]
Chain them: if $\abs{x-a}<\delta$, then $\abs{g(x)-b}<\eta$, and therefore
$\abs{f(g(x))-f(b)}<\eps$. Since $f(b)=f(g(a))$, this says exactly
$\abs{f(g(x))-f(g(a))}<\eps$. The required $\delta$ exists for every $\eps$, so $f\circ g$
is continuous at $a$.
\end{proof}

\begin{focusbox}
Watch the direction of the argument. We pick the tolerance $\eps$ at the \emph{final}
output, pull it back through $f$ to get a tolerance $\eta$ on $f$'s input, then pull
\emph{that} back through $g$ to get a tolerance $\delta$ on the original input. Tolerances
propagate backward, from the end of the pipeline to the start. This backward flow is the
shape of nearly every chain-rule-style argument in calculus. Hold onto it; it returns when
we differentiate compositions.
\end{focusbox}

\begin{example}[Reading off continuity at a glance]
\label{ch08:ex-glance}
The function $h(x)=\sqrt{\,x^2+1\,}$ is continuous on all of $\R$. Here $g(x)=x^2+1$ is a
polynomial (continuous, and always $\ge 1>0$), and $f(u)=\sqrt{u}$ is continuous on
$[0,\infty)$, which safely contains every value of $g$. By
Theorem~\ref{ch08:thm-composition} the composition is continuous everywhere. Similarly
$\sin(e^x)$, $e^{\cos x}$, and $\dfrac{1}{1+x^2}$ are continuous on all of $\R$ --- each is
a finite assembly of continuous atoms by addition, multiplication, division (with nonzero
denominator), and composition. After a little practice you simply \emph{see} it.
\end{example}

\begin{historybox}
For nearly two centuries calculus thrived on this no-jump intuition without a precise
definition. Newton spoke of quantities flowing; Euler manipulated functions as if they
were always smooth. But ``obvious'' continuity hid traps. In the early nineteenth century
Bernard Bolzano (1817) and then Augustin-Louis Cauchy (1821) finally pinned continuity to
limits, roughly as in Definition~\ref{ch08:def-continuity}. Bolzano even gave a correct
proof of the Intermediate Value Theorem --- decades ahead of his time, and largely
ignored. The real shock came later: in 1872 Karl Weierstrass exhibited a function that is
continuous \emph{everywhere} yet differentiable \emph{nowhere} --- an unbroken curve so
jagged it has a corner at every point. That monster proved continuity and smoothness are
genuinely different ideas, and that intuition alone, however good, cannot replace a
definition. The $\eps$--$\delta$ machinery you are using is the tool that finally tamed
the monsters.
\end{historybox}

\section{The Intermediate Value Theorem}
\label{ch08:ivt}

We now reach the first great payoff of continuity. It is a theorem you already believe in
your bones, which is exactly why it is so instructive to see proved.

\begin{intuitionbox}
Suppose a function is continuous on $[a,b]$, starts below some level $N$ at $x=a$, and ends
above $N$ at $x=b$. To get from below $N$ to above $N$ \emph{without jumping}, the graph
must cross the line $y=N$ somewhere in between. A continuous pencil stroke that begins below
a horizontal line and ends above it has to touch that line. That is the entire content of
the Intermediate Value Theorem. The hard part is not believing it --- it is proving that
the word ``continuous'' is strong enough to force the crossing, using only the definition.
\end{intuitionbox}

\begin{theorem}[Intermediate Value Theorem, IVT]
\label{ch08:thm-ivt}
Let $f$ be continuous on the closed interval $[a,b]$, and let $N$ be any number strictly
between $f(a)$ and $f(b)$. Then there exists at least one $c\in(a,b)$ with $f(c)=N$.
\end{theorem}

\begin{figure}[H]
\centering
\begin{tikzpicture}
\begin{axis}[
  width=10cm, height=6cm,
  axis lines=middle,
  xlabel={$x$}, ylabel={$y$},
  xmin=-0.3, xmax=5.3, ymin=-0.5, ymax=4.5,
  xtick={0.5,4.5}, xticklabels={$a$,$b$}, ytick={\empty},
  samples=150, domain=0.5:4.5,
  clip=false,
]
\addplot[curveblue, very thick]{1.0+0.55*(x-0.5)+0.45*sin(deg(0.9*(x-0.5)))};
\draw[targetgreen, dashed, thick] (axis cs:-0.3,2.6)--(axis cs:5.3,2.6);
\node[targetgreen, anchor=west] at (axis cs:4.6,2.6) {\small $y=N$};
\addplot[only marks, mark=*, mark size=2pt, trapred] coordinates {(0.5,1.0)};
\node[trapred,anchor=north] at (axis cs:0.5,0.9){\small $f(a)$};
\addplot[only marks, mark=*, mark size=2pt, trapred] coordinates {(4.5,3.0)};
\node[trapred,anchor=south] at (axis cs:4.4,3.1){\small $f(b)$};
\addplot[only marks, mark=*, mark size=2.5pt, structurepurple] coordinates {(2.643,2.6)};
\node[structurepurple,anchor=north west] at (axis cs:2.7,2.5){\small $c$};
\end{axis}
\end{tikzpicture}
\caption{The Intermediate Value Theorem. A continuous curve running from height $f(a)$ up
to height $f(b)$ must meet every intermediate level $N$ at least once; here it crosses
$y=N$ at the point $c$. For a wigglier curve there could be several crossings, but the
theorem guarantees at least one.}
\label{ch08:fig-ivt}
\end{figure}

\subsection{Why it needs completeness, not just continuity}

Before the proof, a warning that reveals what is really at stake. The IVT is \emph{false}
if we work over the rational numbers alone. Consider $f(x)=x^2-2$ on $[0,2]$, restricted to
rational inputs. It is ``continuous'' in the obvious sense, $f(0)=-2<0$ and $f(2)=2>0$, yet
there is \emph{no rational} $c$ with $f(c)=0$ --- the crossing happens at $\sqrt 2$, which
is not a rational number. The curve sneaks through a gap in the rationals without ever
landing on a rational point.

\begin{happeningbox}
The IVT is not really a theorem about continuity alone. It is a theorem about continuity
\emph{plus the absence of gaps in the real line}. That gaplessness has a name:
\textbf{completeness}, the property that every nonempty set of reals that is bounded above
has a least upper bound (a supremum). The rationals fail completeness --- the set of
rationals whose square is below $2$ has no rational least upper bound --- and the IVT fails
right along with it. The real numbers were, in a sense, \emph{built} to make theorems like
the IVT true: $\R$ is the gapless completion of $\Q$. So when we prove the IVT, the
decisive step will be the one place we invoke a least upper bound. Watch for it.
\end{happeningbox}

\subsection{Proof by ``the last moment it is still below''}

There are two standard proofs. One repeatedly bisects the interval, trapping the crossing
in shrinking subintervals. The other, which we give, captures the crossing in a single
clean stroke using completeness. The idea is irresistible once you see it.

\begin{proof}
It is enough to handle the case $f(a)<N<f(b)$; the other case ($f(a)>N>f(b)$) follows by
applying this one to $-f$ and $-N$. Define
\[
  S=\set{\,x\in[a,b] : f(x)<N\,}.
\]
Think of $S$ as the set of inputs at which the curve is still strictly below the target
level. It is nonempty, since $a\in S$ (because $f(a)<N$), and it is bounded above by $b$.
By the \textbf{completeness} of $\R$, $S$ has a least upper bound; call it $c=\sup S$.
Intuitively, $c$ is \emph{the last moment the curve is still below $N$} --- the exact place
where ``below'' must hand off to ``above.'' We claim $f(c)=N$, and we rule out the only two
alternatives.

\textbf{Suppose $f(c)<N$.} Then $c\ne b$ (since $f(b)>N$), so $c<b$. Because $f$ is
continuous at $c$ and $f(c)$ is strictly less than $N$, values of $f$ stay below $N$ on a
whole little interval to the right of $c$: taking $\eps=N-f(c)>0$, continuity gives a
$\delta>0$ with $f(x)<f(c)+\eps=N$ for all $x$ within $\delta$ of $c$. So some points just
to the \emph{right} of $c$ also lie in $S$. But then $c$ was not an upper bound of $S$ ---
contradiction. The curve cannot still be below $N$ \emph{and} have room to its right.

\textbf{Suppose $f(c)>N$.} Then $c\ne a$, so $c>a$. By the same continuity argument with
$\eps=f(c)-N>0$, the values of $f$ stay \emph{above} $N$ on a whole little interval around
$c$, including a stretch to the \emph{left} of $c$. On that stretch no point belongs to
$S$, so $c-\delta$ would already be an upper bound for $S$ --- smaller than $c$,
contradicting that $c$ is the \emph{least} upper bound.

Both alternatives are impossible, so the only survivor is $f(c)=N$. And $c$ lies strictly
between $a$ and $b$ because $f(a)<N<f(b)$ rules out the endpoints.
\end{proof}

\begin{focusbox}
Feel where each hypothesis earned its keep. \emph{Completeness} produced the candidate
point $c=\sup S$ --- without it, $S$ might have no least upper bound and the proof could not
even begin. \emph{Continuity} was used twice, once on each side, to show that the curve
cannot be strictly below or strictly above $N$ at $c$ without contradicting how $c$ was
chosen. Drop continuity and a jump could leap over $N$; drop completeness and the crossing
could fall into a gap. The theorem lives exactly at the meeting of these two ideas.
\end{focusbox}

\subsection{Root-finding: the IVT as an existence machine}

The IVT is an \emph{existence} theorem. It does not tell you \emph{where} the crossing is,
only that one exists. That modest-sounding guarantee is the foundation of every numerical
root-finder.

\begin{corollary}[Bolzano's sign-change test]
\label{ch08:cor-rootsign}
If $f$ is continuous on $[a,b]$ and $f(a)$ and $f(b)$ have opposite signs (one negative,
one positive), then $f$ has a root in $(a,b)$: some $c$ with $f(c)=0$.
\end{corollary}

This is just the IVT with $N=0$: the value $0$ lies between a negative number and a
positive number, so it is attained.

\begin{example}[A root of a quintic you cannot solve by formula]
\label{ch08:ex-quintic}
Does $p(x)=x^5+x-1$ have a real root? There is no algebraic formula for the roots of a
general quintic, but we do not need one. $p$ is a polynomial, hence continuous everywhere.
Compute the ends:
\[
  p(0)=0+0-1=-1<0, \qquad p(1)=1+1-1=1>0.
\]
Opposite signs across $[0,1]$, so by Corollary~\ref{ch08:cor-rootsign} there is a root in
$(0,1)$ (Figure~\ref{ch08:fig-quintic}). Existence, established in two arithmetic lines and
one theorem, with no formula in sight.
\end{example}

\begin{figure}[H]
\centering
\begin{tikzpicture}
\begin{axis}[
  width=10cm, height=5.6cm,
  axis lines=middle,
  xlabel={$x$}, ylabel={$y$},
  xmin=-0.1, xmax=1.15, ymin=-1.3, ymax=1.3,
  xtick={0,0.5,0.755,1}, xticklabels={$0$,$0.5$,$c$,$1$}, ytick={-1,1},
  samples=120, domain=0:1.1,
  clip=false,
]
\addplot[curveblue, very thick]{x^5+x-1};
\addplot[only marks, mark=*, mark size=2pt, trapred] coordinates {(0,-1)};
\node[trapred, anchor=north east] at (axis cs:0,-1){\small $p(0)=-1$};
\addplot[only marks, mark=*, mark size=2pt, targetgreen] coordinates {(1,1)};
\node[targetgreen, anchor=south east] at (axis cs:1,1){\small $p(1)=1$};
\addplot[only marks, mark=*, mark size=2.5pt, structurepurple] coordinates {(0.7549,0)};
\node[structurepurple, anchor=south west] at (axis cs:0.78,0.05){\small root $c$};
\end{axis}
\end{tikzpicture}
\caption{The quintic $p(x)=x^5+x-1$ on $[0,1]$. Since $p(0)=-1<0$ and $p(1)=1>0$, the
continuous curve must cross the axis; the IVT guarantees a root $c$ (here $c\approx0.755$)
even though no algebraic formula produces it.}
\label{ch08:fig-quintic}
\end{figure}

\begin{labbox}
Pin that quintic root down by hand using \emph{bisection} --- the IVT applied over and
over. We know the root lies in $[0,1]$ since $p(0)<0<p(1)$. Test the midpoint:
\[
  p(0.5)=0.5^5+0.5-1=0.03125-0.5=-0.46875<0.
\]
Same sign as $p(0)$, so the root is in $[0.5,1]$. Now $p(0.75)=0.2373+0.75-1=-0.0127<0$, so
the root is in $[0.75,1]$. Next $p(0.875)\approx 0.5129+0.875-1=0.388>0$, so the root is in
$[0.75,0.875]$. Each step halves the interval and the IVT promises the root never escapes.
Two more halvings put you near $x\approx 0.755$. You have computed a root of an unsolvable
quintic with nothing but arithmetic and a theorem.
\end{labbox}

\begin{whybox}
Why is bisection guaranteed to work, and not just lucky? Each step keeps the property ``the
endpoints have opposite signs'' as an \emph{invariant}. As long as that invariant holds, the
IVT guarantees a root inside the current interval. Halving the interval shrinks its length
to $(b-a)/2^n$ after $n$ steps, which tends to $0$, so the trapped root is pinned to
arbitrary precision. The method cannot fail to converge, because at every stage a root is
provably present in a provably shrinking cage. Bisection is the IVT turned into an
algorithm.
\end{whybox}

\subsection{Fixed points: where a function meets the diagonal}

A second beautiful use. A \textbf{fixed point} of $g$ is an input that $g$ leaves
unchanged: $g(c)=c$. Geometrically, it is where the graph of $g$ meets the diagonal line
$y=x$.

\begin{corollary}[A fixed-point theorem]
\label{ch08:cor-fixed}
If $g$ is continuous and maps $[a,b]$ into itself --- that is, $a\le g(x)\le b$ for all
$x\in[a,b]$ --- then $g$ has a fixed point in $[a,b]$.
\end{corollary}

\begin{proof}
Consider the auxiliary function $h(x)=g(x)-x$, which is continuous as a difference of
continuous functions. At the left end, $g(a)\ge a$ gives $h(a)=g(a)-a\ge 0$. At the right
end, $g(b)\le b$ gives $h(b)=g(b)-b\le 0$. So $h$ runs from a nonnegative value down to a
nonpositive one. If either end is exactly $0$ we have our fixed point outright; otherwise
$h(a)>0>h(b)$, and the IVT hands us a $c$ with $h(c)=0$, i.e.\ $g(c)=c$.
\end{proof}

\begin{figure}[H]
\centering
\begin{tikzpicture}
\begin{axis}[
  width=7.6cm, height=6cm,
  axis lines=middle,
  xlabel={$x$}, ylabel={$y$},
  xmin=-0.2, xmax=2.4, ymin=-0.2, ymax=2.4,
  xtick={0,2}, xticklabels={$a$,$b$}, ytick={\empty},
  samples=120, domain=0:2,
  clip=false,
]
\addplot[gray, thick, domain=-0.2:2.4]{x};
\node[gray, anchor=south east] at (axis cs:2.3,2.3){\small $y=x$};
\addplot[curveblue, very thick]{0.4+1.1*sin(deg(0.8*x))};
\node[curveblue, anchor=north west] at (axis cs:0.05,1.1){\small $y=g(x)$};
\addplot[only marks, mark=*, mark size=2.5pt, structurepurple] coordinates {(1.384,1.384)};
\node[structurepurple, anchor=north west] at (axis cs:1.45,1.3){\small fixed point};
\end{axis}
\end{tikzpicture}
\caption{A fixed point is a crossing of $y=g(x)$ with the diagonal $y=x$. If a continuous
$g$ keeps its outputs inside $[a,b]$, its graph starts on or above the diagonal at $a$ and
ends on or below it at $b$, so it must cross.}
\label{ch08:fig-fixed}
\end{figure}

\begin{connectionbox}
This one-dimensional fixed-point theorem is the humble ancestor of a towering family. In
higher dimensions it becomes Brouwer's Fixed-Point Theorem (every continuous map of a disk
to itself has a fixed point), which underlies the existence of economic equilibria
(Nash's theorem on equilibria leans on it) and of solutions to differential equations. The
everyday image: stir a cup of coffee continuously, and when it settles, at least one
molecule sits exactly where it started. Stretch and fold a paper map of a region and lay it
back inside the same region --- some point lands on top of where it began. All of it grows
from the same seed: a continuous self-map cannot avoid the diagonal.
\end{connectionbox}

\section{The Extreme Value Theorem}
\label{ch08:evt}

The IVT guarantees a function hits intermediate \emph{values}. Our second great theorem
guarantees it hits extreme ones --- a genuine highest point and a genuine lowest point ---
provided we ask on a closed, bounded interval.

\begin{theorem}[Extreme Value Theorem, EVT]
\label{ch08:thm-evt}
If $f$ is continuous on a closed bounded interval $[a,b]$, then $f$ attains a maximum and a
minimum there: there exist points $p,q\in[a,b]$ with
\[
  f(q)\le f(x)\le f(p)\qquad\text{for all } x\in[a,b].
\]
\end{theorem}

The word \emph{attains} is the heart of it. It is not merely that the values stay bounded
--- it is that the largest value is actually \emph{achieved} by some input $p$, and the
smallest by some input $q$. The maximum is a value the function truly reaches, not just a
ceiling it approaches.

\subsection{Why every hypothesis is necessary}

This theorem is fragile: weaken any single condition and it collapses. Studying the
failures shows exactly what each word is buying.

\begin{itemize}
  \item \textbf{Drop ``closed'' (omit an endpoint).} On the \emph{open} interval $(0,1)$,
        $f(x)=x$ is continuous but attains neither a max nor a min: it gets arbitrarily
        close to $1$ and to $0$ without ever reaching them. The supremum $1$ and infimum
        $0$ exist as bounds but are never \emph{achieved}, because the points that would
        achieve them were excluded from the domain.
  \item \textbf{Drop ``bounded'' (let the interval run to infinity).} On $[0,\infty)$,
        $f(x)=x$ is continuous but has no maximum --- it grows without bound. A maximum
        cannot exist where the values escape to infinity.
  \item \textbf{Drop ``continuous.''} On $[0,1]$, let $f(x)=1/x$ for $x>0$ and $f(0)=0$.
        This is bounded below but shoots to $+\infty$ near $0$, so it has no maximum despite
        living on a perfectly closed bounded interval. A single infinite discontinuity is
        enough to destroy the conclusion.
\end{itemize}

\begin{figure}[H]
\centering
\begin{tikzpicture}
\begin{axis}[
  width=10.5cm, height=5.6cm,
  axis lines=middle,
  xlabel={$x$}, ylabel={$y$},
  xmin=-0.3, xmax=5.3, ymin=-0.4, ymax=3.6,
  xtick={0.4,4.6}, xticklabels={$a$,$b$}, ytick={\empty},
  samples=150, domain=0.4:4.6,
  clip=false,
]
\addplot[curveblue, very thick]{1.7+1.2*sin(deg(1.4*(x-0.4)))};
\addplot[only marks, mark=*, mark size=2.5pt, trapred] coordinates {(1.52,2.9)};
\node[trapred, anchor=south] at (axis cs:1.52,3.0){\small max $f(p)$};
\addplot[only marks, mark=*, mark size=2.5pt, targetgreen] coordinates {(3.76,0.5)};
\node[targetgreen, anchor=north] at (axis cs:3.76,0.4){\small min $f(q)$};
\draw[trapred, dashed] (axis cs:-0.3,2.9)--(axis cs:1.52,2.9);
\draw[targetgreen, dashed] (axis cs:-0.3,0.5)--(axis cs:3.76,0.5);
\end{axis}
\end{tikzpicture}
\caption{The Extreme Value Theorem. A continuous function on a closed bounded interval
genuinely reaches a highest value $f(p)$ and a lowest value $f(q)$. The dashed lines mark
the ceiling and floor, each \emph{touched} by the curve.}
\label{ch08:fig-evt}
\end{figure}

\begin{focusbox}
The EVT is the silent partner behind all of optimization. Later, when we hunt for maxima
and minima by setting derivatives to zero, we will be searching among \emph{candidate}
points --- critical points and endpoints. But why should a maximum exist among them at all?
Because the EVT promises one exists \emph{somewhere}, and a short argument then shows that
``somewhere'' has to be a critical point or an endpoint. Without the EVT, we would be
hunting for something that might not be there. It guarantees the prize is on the table
before we start looking for it.
\end{focusbox}

\subsection{The shape of the argument}

A complete proof leans on two facts about continuous functions on $[a,b]$ that are slightly
beyond our scope to prove in full, but whose ideas we can convey honestly. We sketch the
route so the theorem does not feel like magic.

\begin{intuitionbox}
\textbf{Step 1: the function is bounded.} First show $f$ cannot run off to $\pm\infty$ on
$[a,b]$. The mechanism: if $f$ were unbounded, you could find inputs $x_1,x_2,\dots$ in
$[a,b]$ with $\abs{f(x_n)}\to\infty$. Because $[a,b]$ is closed and bounded, these inputs
must \emph{accumulate} somewhere inside it --- they cannot all spread out to infinity (the
interval is finite) and cannot escape out an open end (there is none). Near that
accumulation point, continuity forces the $f$-values to stay close to one finite number,
flatly contradicting $\abs{f(x_n)}\to\infty$. So $f$ must be bounded.

\textbf{Step 2: the bound is achieved.} Let $M$ be the least upper bound of the values
$f(x)$ --- it exists by completeness now that we know the values are bounded. By definition
of least upper bound, there are inputs with $f$-values as close to $M$ as we like; collect
such inputs $x_1,x_2,\dots$ with $f(x_n)\to M$. Again they accumulate at some point $p$ in
$[a,b]$, and continuity at $p$ forces $f(p)=M$. So the supremum $M$ is not just a ceiling ---
it is actually \emph{attained} at $p$. Apply the same reasoning to $-f$ to attain the
minimum.
\end{intuitionbox}

\begin{happeningbox}
Both steps secretly use the same precious property of $[a,b]$: \emph{any infinite sequence
of points in a closed bounded interval has a point it crowds up against, inside the
interval.} This is the Bolzano--Weierstrass property, and it is what ``closed and bounded''
really buys you. ``Bounded'' keeps the points from escaping to infinity; ``closed'' keeps
the accumulation point from leaking out through a missing endpoint. Strip away either and
the accumulation point can vanish --- which is precisely why the counterexamples on $(0,1)$
and $[0,\infty)$ work. Continuity then does the final job: it transfers the accumulation of
\emph{inputs} into a statement about \emph{outputs}, letting a limiting value be genuinely
reached. The EVT is the marriage of compactness (closed + bounded) with continuity.
\end{happeningbox}

\begin{connectionbox}
The phrase ``closed and bounded'' is the one-dimensional face of a concept that organizes
all of analysis and topology: \textbf{compactness}. The grand statement, of which the EVT
is a special case, is that \emph{a continuous real-valued function on a compact set attains
its maximum and minimum}. Compactness is exactly the abstraction of ``no points can escape''
that we needed in both steps above. When you later meet compactness in a topology course,
you can think back to this chapter: it is the property that makes the Extreme Value Theorem
true, distilled to its essence and freed from the real line.
\end{connectionbox}

\section{Why continuity is the right notion of ``no jumps''}
\label{ch08:right-notion}

We close by stepping back. We \emph{defined} continuity by a limit equation, but is that
definition the right one --- does it really capture ``no jumps,'' nothing more and nothing
less? The two big theorems are the evidence.

A function with a jump fails the IVT: the Heaviside step $H$ leaps from $0$ to $1$ at the
origin and skips every value strictly between, like $\tfrac12$. So ``jumpy'' functions are
exactly the ones that can dodge intermediate values, and continuous functions are exactly
the ones that cannot. The IVT is, in effect, the theorem that says \emph{our definition of
continuity is strong enough to forbid jumps.} Conversely, our gallery showed that every way
a function can fail Definition~\ref{ch08:def-continuity} --- hole, jump, pole, oscillation
--- is visibly a break in the curve. The definition forbids breaks, and only breaks.

\begin{aestheticbox}
There is a deep economy here worth pausing on. We wanted a geometric idea --- an unbroken
curve --- and we found we could capture it with a single algebraic equation,
$\lim_{x\to a}f(x)=f(a)$, that mentions no pictures at all. Then that one humble equation,
applied with the completeness of $\R$, forced two sweeping existence theorems: every
intermediate value is hit (IVT), and extreme values are genuinely attained (EVT). This is
the recurring miracle of analysis: a precise local condition (how a function behaves near
each single point) propagates into powerful global conclusions (what the function must do
across an entire interval). Local restraint, global consequence. The whole subject runs on
this current, and continuity is where you first feel it.
\end{aestheticbox}

\begin{recoverybox}
Before the exercises, the through-line in one breath. Continuity at $a$ means the limit
equals the value, which is the precise form of ``no jump.'' Continuity is inherited through
sums, products, quotients (away from zeros), and compositions, so almost every function you
write down is continuous on its natural domain. On the real line --- gapless by completeness
--- continuity forces the Intermediate Value Theorem (every in-between value is attained,
the engine of root-finding and fixed points) and, on closed bounded intervals, the Extreme
Value Theorem (max and min are truly reached, the foundation of optimization). Every one of
these rests on the same modest equation. If a step ever feels mysterious, return to the
picture of a pencil that never leaves the page.
\end{recoverybox}

\section{Exercises}
\label{ch08:exercises}

\begin{exercise}[Classify the discontinuity]
\label{ch08:exr-classify}
For each function, find every point of discontinuity on $\R$ and classify it as removable,
jump, infinite, or oscillatory.
\[
  \text{(a) } \frac{x^2-9}{x-3}, \quad
  \text{(b) } \frac{1}{(x-2)^2}, \quad
  \text{(c) } \frac{\abs{x}}{x}, \quad
  \text{(d) } \cos\!\Big(\frac{1}{x}\Big).
\]
For any removable case, state the value that repairs it.
\end{exercise}

\begin{exercise}[Make it continuous]
\label{ch08:exr-piecewise}
Find the value of $k$ that makes
\[
  f(x)=\begin{cases} \dfrac{x^2-4}{x-2}, & x\ne 2,\\[1.4ex] k, & x=2,\end{cases}
\]
continuous at $x=2$. Then do the same for
$\displaystyle g(x)=\dfrac{\sqrt{x+1}-1}{x}$ at $x=0$ (rationalize the numerator first).
\end{exercise}

\begin{exercise}[Continuity by inheritance]
\label{ch08:exr-inheritance}
Without using $\eps$ and $\delta$, explain why each function is continuous on the stated
set, naming which combination or composition laws you use:
(a) $\dfrac{\sin x}{x^2+1}$ on $\R$; \quad
(b) $\sqrt{\,4-x^2\,}$ on $[-2,2]$; \quad
(c) $e^{1/x}$ on $(0,\infty)$.
\end{exercise}

\begin{exercise}[A root exists]
\label{ch08:exr-root}
Show that the equation $\cos x = x$ has a solution in $\big[0,\tfrac{\pi}{2}\big]$. (Hint:
apply the IVT to $h(x)=\cos x - x$ and check the endpoints.) Then use one bisection step to
locate the solution in a smaller interval.
\end{exercise}

\begin{exercise}[A fixed point in the wild]
\label{ch08:exr-fixed}
Let $g(x)=\tfrac12\cos x$. Show $g$ maps the interval $[-1,1]$ into itself and is
continuous, hence has a fixed point there. What equation does the fixed point satisfy?
\end{exercise}

\begin{exercise}[An interval always has a midpoint value]
\label{ch08:exr-avg}
Suppose $f$ is continuous on $[a,b]$. Prove there is a point $c\in[a,b]$ with
\[
  f(c)=\frac{f(a)+f(b)}{2}.
\]
(The average of the two endpoint heights lies between them.) Is the same true for
$\frac{f(a)+2f(b)}{3}$? Explain.
\end{exercise}

\begin{exercise}[Every hypothesis matters]
\label{ch08:exr-hyp}
Give an explicit example of a function that is continuous on the half-open interval
$(0,1]$ and is \emph{unbounded} there, showing the EVT genuinely needs the left endpoint.
Then give a function continuous and bounded on $(0,1]$ that nonetheless attains no maximum.
\end{exercise}

\begin{exercise}[The temperature on a great circle, a proof]
\label{ch08:exr-antipodal}
Temperature varies continuously along the equator. Prove there exist two diametrically
opposite points on the equator with exactly the same temperature. Model the equator by
$[0,2\pi]$ with $T(0)=T(2\pi)$, and apply the IVT to the function
$D(\theta)=T(\theta+\pi)-T(\theta)$ on $[0,\pi]$. (Hint: examine $D(0)$ and $D(\pi)$ and
notice they have opposite signs unless one is already zero.)
\end{exercise}

\begin{exercise}[A continuous bijection's inverse, conceptual]
\label{ch08:exr-monotone}
Argue, using the IVT, that a continuous function on an interval that is one-to-one must be
strictly monotonic (always increasing or always decreasing). (Hint: if it rose then fell,
the IVT would force it to repeat a value, contradicting one-to-oneness.) This is the secret
reason continuous invertible functions have continuous inverses --- a fact we will use when
differentiating $\ln x$ and $\arcsin x$.
\end{exercise}

\begin{exercise}[Prove a combination law]
\label{ch08:exr-prove}
Using only Definition~\ref{ch08:def-continuity} and the limit laws, prove carefully that if
$f$ and $g$ are continuous at $a$ then $fg$ is continuous at $a$. Then explain in one
sentence why your proof breaks down for $f/g$ when $g(a)=0$.
\end{exercise}

\begin{exercise}[Dirichlet's monster, a challenge]
\label{ch08:exr-dirichlet}
Let $D(x)=1$ if $x$ is rational and $D(x)=0$ if $x$ is irrational. Show $D$ is discontinuous
at \emph{every} real number. (Hint: in any interval, however small, there are both rationals
and irrationals, so the values $0$ and $1$ both occur arbitrarily close to any point, and no
single limit can exist.) This function has no points of continuity at all --- the opposite
extreme from the well-behaved functions of this chapter.
\end{exercise}
